\section{GlyphChain two-phase production-release authorization}
\label{sec:glyphchain-two-phase-release-authorization}

The final production decision is deliberately separated from readiness evidence.
Passing technical, operational, security and Legal gates demonstrates that a
specific release candidate is eligible for consideration; it does not itself
authorize deployment.  Tessaris therefore implements a two-phase protocol in
which the first thirty-three gates are frozen into a preauthorization report and
the thirty-fourth gate signs an exact commitment to that report.

\subsection{Phase one: clean preauthorization state}

Let $T$ be the canonical set of thirty-four release tasks and let
$t_f=\texttt{REL-10}$ be the final production-authorization task.  For release
manifest digest $r$, authenticated trust-root inventory $G$ and evidence inventory
$E$, the evaluator may emit
\texttt{AWAITING\_FINAL\_RELEASE\_AUTHORIZATION} only when

\begin{equation}
  \begin{split}
  P_{33}(r,E,G) ={}&
    \bigwedge_{t\in T\setminus\{t_f\}}\operatorname{pass}(t)
    \;\land\;\neg\operatorname{pass}(t_f) \\
    &\land\;\operatorname{absentEvidence}(t_f)
    \land\;\operatorname{authenticated}(G) \\
    &\land\;\neg\operatorname{unknownTask}(E)
    \land\;\neg\operatorname{duplicateTask}(E)
    \land\;\neg\operatorname{globalError}.
  \end{split}
\end{equation}

The final gate must be absent rather than merely malformed or rejected.  Its task
record has no evidence-body digest and no attestation identifiers.  This prevents
an attempted early authorization from being relabelled as a clean waiting state.
The report records the exact release-manifest digest, evidence-inventory digest,
trust-root inventory file digest, governance-authority key identifier, every
per-task result and its own canonical report digest.

Before a phase-one file can be used by phase two, the core evaluator
revalidates the complete report structure.  It requires exactly thirty-three
passing task identifiers, exactly one unmet identifier
\texttt{REL-10}, authenticated and structurally valid trust roots, no unknown or
duplicate tasks, no global errors, and a passing result for every individual
preceding task.  It also verifies that the report binds the same release manifest,
trust-root file, governance key and trust-root map as the proposed final
evaluation.  The phase-one evidence-inventory digest must equal the digest of the
final inventory after removing only the \texttt{REL-10} entry.  Consequently,
adding final authorization is permitted, but silently changing any preceding
entry---even an otherwise ignored top-level field---invalidates the phase-one
continuity check.

\subsection{Phase-two evidence contract}

The final \texttt{REL-10} evidence body requires both of the following retained
artefact classes:

\begin{itemize}
  \item \texttt{preauthorization\_readiness\_report}, containing the exact
        phase-one readiness-report bytes; and
  \item \texttt{production\_release\_authorization}, containing the explicit
        controlled decision to release that candidate.
\end{itemize}

If $F_{33}$ denotes the exact phase-one report file, the signed artefact commitment
must contain

\begin{equation}
  h_{33}=\operatorname{SHA256}(F_{33}).
\end{equation}

The final evaluator reads the report once as a bounded regular file, refuses a
symbolic link, validates its internal digest and state, and computes $h_{33}$ from
the same byte sequence.  The \texttt{REL-10} artefact set must contain exactly that
digest for the preauthorization kind.  An omitted report produces
\texttt{preauthorization\_report\_required}; an unvalidated report produces
\texttt{preauthorization\_report\_not\_validated}; and a different digest produces
\texttt{preauthorization\_report\_binding\_invalid}.

This binding closes a time-of-check/time-of-authorization gap.  The final
authorities cannot sign a generic statement and later attach a different readiness
result, nor can evidence from one manifest, preceding evidence inventory or
authority set authorize another.  The core evaluator accepts the bounded raw
report bytes, parses and validates them itself, and derives the file digest; callers
cannot supply a trusted boolean declaring that validation already happened.

\subsection{Final authority and separation of duties}

The final body is signed over the domain-separated canonical release-evidence
payload by three distinct configured authorities:

\begin{enumerate}
  \item the release approver;
  \item the Legal product owner; and
  \item the application-security reviewer.
\end{enumerate}

The evaluator performs an injective role assignment, so one signer cannot satisfy
two required roles.  At least one assigned signer must be marked independent in the
governance-authenticated trust-root inventory.  Duplicate public keys under
different signer names are rejected across the entire root inventory, preventing
cosmetic identities from bypassing separation of duties.

Let $S_f$ represent valid distinct final signatures, $A_f$ the two required
artefact classes, $B_f$ the exact $h_{33}$ binding and $V_{33}$ successful
revalidation of the preauthorization report.  Production authorization is

\begin{equation}
  \operatorname{Authorize}(r) =
  P_{33}(r,E,G) \land V_{33} \land B_f \land A_f \land S_f.
\end{equation}

No percentage, user-interface toggle or presence of a filename substitutes for
this conjunction.  Only a final report with all thirty-four gates passing emits
\texttt{ALL\_RELEASE\_GATES\_EVIDENCED} and sets
\texttt{production\_release\_authorized=true}.

\subsection{Operational workflow}

The controlled workflow is:

\begin{enumerate}
  \item freeze the candidate release manifest;
  \item assemble and independently countersign evidence for gates
        \texttt{LD-01} through \texttt{REL-09};
  \item evaluate that inventory and retain the resulting clean phase-one report;
  \item present that immutable report and its exact digest to the three final
        authorities;
  \item create the \texttt{REL-10} body with both required artefacts and the exact
        preauthorization-report digest;
  \item obtain the three distinct Ed25519 attestations without sharing private
        keys; and
  \item rerun the evaluator with the phase-one report supplied explicitly through
        \texttt{--preauthorization-report}.
\end{enumerate}

Any change to the release manifest, evidence inventory, trust-root inventory or
governance anchor after phase one requires a new phase-one evaluation and new final
authorization.  A former valid authorization is audit evidence for its historical
candidate, not permission for a modified release.

\subsection{Implemented verification and claim boundary}

Regression tests exercise the blocked state, the exact thirty-three-gate waiting
state, successful transition to all thirty-four gates, omitted preauthorization,
wrong report digest, modification of a preceding evidence entry, distinct-role
assignment and authenticated trust-root requirements.  The evaluation report
preserves the derived preauthorization file digest, validation result and explicit
validation errors so downstream audit tooling can distinguish a mere byte sequence
from a checked phase-one report.

This protocol proves that configured authorities explicitly approved the exact
candidate after the configured evidence gates passed.  It does not prove that an
external report is truthful merely because its bytes were signed, replace qualified
professional review, create qualified electronic-signature status, or authorize a
release when any external gate remains unmet.  Those substantive facts remain the
responsibility of the named reviewers and the retained external artefacts.
